% !PRINTED 121
\weberpage{121}
%<<CONT
und die Funktion $\sigma(u)$ selbst kann entsprechend auch mit $\sigma_{11}(u)$ bezeichnet werden.

Für diese Funktionen ergeben sich nach (4) die charakteristischen Periodengleichungen

\weq{6}{\begin{aligned}
\sigma_{g_1,g_2}(u+\omega_1) &= (-1)^{g_1}\,e^{\eta_1(2u+\omega_1)}\,\sigma_{g_1,g_2}(u),\\
\sigma_{g_1,g_2}(u+\omega_2) &= (-1)^{g_2}\,e^{\eta_2(2u+\omega_2)}\,\sigma_{g_1,g_2}(u),\\
\sigma_{g_1,g_2}(u+a\omega_1+b\omega_2) &= (-1)^{g_1 a+g_2 b+ab}\,e^{(a\eta_1+b\eta_2)(2u+a\omega_1+b\omega_2)}\,\sigma_{g_1,g_2}(u).
\end{aligned}}

Durch Anwendung einer linearen Transformation werden die drei Funktionen $\sigma_{00}, \sigma_{01}, \sigma_{10}$ untereinander permutiert, wie die Formel (2) lehrt [oder \S\,35, (1)]. Je nach dieser Permutation zerfallen die linearen Transformationen in \gesperrt{sechs Klassen}, deren erste alle die Transformationen umfaßt, die die Funktionen $\sigma_{00}, \sigma_{01}, \sigma_{10}$ ungeändert lassen. Diese sind dadurch charakterisiert, daß $\alpha, \delta$ ungerade, $\beta, \gamma$ gerade Zahlen sind, was wir kurz so schreiben:

\[
\begin{pmatrix}\alpha, & \beta\\ \gamma, & \delta\end{pmatrix} \equiv \begin{pmatrix}1, & 0\\ 0, & 1\end{pmatrix} (\mathrm{mod}\ 2).
\]

Hiernach sind die sechs Klassen der linearen Transformationen folgendermaßen zu charakterisieren:

\weq{7}{\begin{aligned}
&\text{I.}\quad \begin{pmatrix}\alpha, & \beta\\ \gamma, & \delta\end{pmatrix} \equiv \begin{pmatrix}1, & 0\\ 0, & 1\end{pmatrix} (\mathrm{mod}\ 2)\quad (00, 01, 10),\\
&\text{II.}\quad \text{,,} \equiv \begin{pmatrix}0, & 1\\ -1, & 0\end{pmatrix} \quad \text{,,}\quad (00, 10, 01),\\
&\text{III.}\quad \text{,,} \equiv \begin{pmatrix}1, & 1\\ -1, & 0\end{pmatrix} \quad \text{,,}\quad (10, 00, 01),\\
&\text{IV.}\quad \text{,,} \equiv \begin{pmatrix}1, & 1\\ 0, & 1\end{pmatrix} \quad \text{,,}\quad (10, 01, 00),\\
&\text{V.}\quad \text{,,} \equiv \begin{pmatrix}1, & 0\\ 1, & 1\end{pmatrix} \quad \text{,,}\quad (01, 00, 10),\\
&\text{VI.}\quad \text{,,} \equiv \begin{pmatrix}0, & 1\\ -1, & 1\end{pmatrix} \quad \text{,,}\quad (01, 10, 00),
\end{aligned}}
wo in der letzten Kolumne die jedesmalige Permutation der Charakteristiken 00, 01, 10 aufgeführt ist.

Die Transformationen der ersten Klasse bilden eine in der Gruppe \gesperrt{aller} linearen Substitutionen enthaltene Gruppe $\mathfrak{A}$, also einen Teiler der Gruppe $\mathfrak{L}$ (\S\,28). Setzen wir
%CONT>>
